% GENERATED FILE. Edit canonical lessons, then run npm run generate:books.
% canonical-source-hash: fnv1a64:c7f749109658efe6
% canonical-lessons: IT-C443-settentrionale, IT-C443-meridionale, IT-C443-orientale, IT-C443-occidentale, IT-C443-punto-di-riferimento

\chapter{I punti cardinali}
\label{ch:it-i-punti-cardinali}
\hlchaptermodality{443}

\begin{chapteropening}
\textbf{By the end of this chapter:} \emph{I can place things on the compass: northern, southern, eastern, western, and a landmark.}

Complete the last lesson of chapter 443: I can place things on the compass: northern, southern, eastern, western, and a landmark.
\end{chapteropening}

\section[settentrionale]{settentrionale — northern}

\label{lesson:IT-C443-settentrionale}



\begin{quote}
Before the new one: say the Italian for to come out, then the Italian for to go round.
\end{quote}

\subsection*{You'll want to know: settentrionale}
\textbf{settentrionale} — \textquotedblleft{}northern\textquotedblright{}.

\begin{quote}
Le regioni settentrionali sono più fredde.
\end{quote}

\begin{quote}
\emph{The northern regions are colder.}
\end{quote}

Formal for \emph{del nord}; \emph{il settentrione} — the north.

\begin{cousinweb}
\emph{settentrionale} is a learned borrowing from Latin \emph{septentriōnālis} 'northern' (EtymDB), from \emph{septentriōnēs} 'the seven plough oxen', the seven stars of the north.
\end{cousinweb}

\begin{grammarlens}[title={What you've built}]
One more word for telling, explaining and describing at B1.
\end{grammarlens}

\subsection*{Your turn}
\begin{itemize}
\raggedright
  \item \emph{Say it:} \emph{settentrionale}
  \item \emph{Say it:} \emph{settentrionale}, once more
  \item \emph{Say it:} a sentence of your own with \emph{settentrionale}
  \item \emph{Recall:} say the Italian for to come out, then the Italian for to go round, then say \emph{settentrionale} again
\end{itemize}

\begin{tcolorbox}[breakable,colback=teal!4,colframe=teal!35!black,title={Before you move on}]
What does settentrionale mean? (\textquotedblleft{}Northern\textquotedblright{}.) Say it once more.
\end{tcolorbox}
\section[meridionale]{meridionale — southern}

\label{lesson:IT-C443-meridionale}



\begin{quote}
Before the new one: say the Italian for to go round, then the Italian for northern.
\end{quote}

\subsection*{You'll want to know: meridionale}
\textbf{meridionale} — \textquotedblleft{}southern\textquotedblright{}.

\begin{quote}
La cucina meridionale ha molte verdure.
\end{quote}

\begin{quote}
\emph{Southern cooking has a lot of vegetables.}
\end{quote}

\emph{Il Meridione} is the South; \emph{un meridionale} is a southerner.

\begin{cousinweb}
\emph{meridionale} is a learned borrowing from Latin \emph{merīdiōnālis} (EtymDB), from \emph{merīdiēs} 'midday': at noon the sun stands in the south.
\end{cousinweb}

\begin{grammarlens}[title={What you've built}]
One more word for telling, explaining and describing at B1.
\end{grammarlens}

\subsection*{Your turn}
\begin{itemize}
\raggedright
  \item \emph{Say it:} \emph{meridionale}
  \item \emph{Say it:} \emph{meridionale}, once more
  \item \emph{Say it:} a sentence of your own with \emph{meridionale}
  \item \emph{Recall:} say the Italian for to go round, then the Italian for northern, then say \emph{meridionale} again
\end{itemize}

\begin{tcolorbox}[breakable,colback=teal!4,colframe=teal!35!black,title={Before you move on}]
What does meridionale mean? (\textquotedblleft{}Southern\textquotedblright{}.) Say it once more.
\end{tcolorbox}
\section[orientale]{orientale — eastern}

\label{lesson:IT-C443-orientale}



\begin{quote}
Before the new one: say the Italian for northern, then the Italian for southern.
\end{quote}

\subsection*{You'll want to know: orientale}
\textbf{orientale} — \textquotedblleft{}eastern\textquotedblright{}.

\begin{quote}
La costa orientale ha il mare più calmo.
\end{quote}

\begin{quote}
\emph{The eastern coast has the calmer sea.}
\end{quote}

\emph{L'arte orientale} — Eastern art.

\begin{cousinweb}
\emph{orientale} is a learned borrowing from Latin \emph{orientālis} 'of the east' (EtymDB), from \emph{oriēns} 'rising (sun)'.
\end{cousinweb}

\begin{grammarlens}[title={What you've built}]
One more word for telling, explaining and describing at B1.
\end{grammarlens}

\subsection*{Your turn}
\begin{itemize}
\raggedright
  \item \emph{Say it:} \emph{orientale}
  \item \emph{Say it:} \emph{orientale}, once more
  \item \emph{Say it:} a sentence of your own with \emph{orientale}
  \item \emph{Recall:} say the Italian for northern, then the Italian for southern, then say \emph{orientale} again
\end{itemize}

\begin{tcolorbox}[breakable,colback=teal!4,colframe=teal!35!black,title={Before you move on}]
What does orientale mean? (\textquotedblleft{}Eastern\textquotedblright{}.) Say it once more.
\end{tcolorbox}
\section[occidentale]{occidentale — western}

\label{lesson:IT-C443-occidentale}



\begin{quote}
Before the new one: say the Italian for southern, then the Italian for eastern.
\end{quote}

\subsection*{You'll want to know: occidentale}
\textbf{occidentale} — \textquotedblleft{}western\textquotedblright{}.

\begin{quote}
Il sole tramonta sul lato occidentale.
\end{quote}

\begin{quote}
\emph{The sun sets on the western side.}
\end{quote}

Built on \emph{occidente}, where the sun goes down.

\begin{cousinweb}
\emph{occidentale} is a learned borrowing from Latin \emph{occidentālis} (EtymDB), from \emph{occidēns} 'setting (sun)', built on \emph{cadere} 'to fall'.
\end{cousinweb}

\begin{grammarlens}[title={What you've built}]
One more word for telling, explaining and describing at B1.
\end{grammarlens}

\subsection*{Your turn}
\begin{itemize}
\raggedright
  \item \emph{Say it:} \emph{occidentale}
  \item \emph{Say it:} \emph{occidentale}, once more
  \item \emph{Say it:} a sentence of your own with \emph{occidentale}
  \item \emph{Recall:} say the Italian for southern, then the Italian for eastern, then say \emph{occidentale} again
\end{itemize}

\begin{tcolorbox}[breakable,colback=teal!4,colframe=teal!35!black,title={Before you move on}]
What does occidentale mean? (\textquotedblleft{}Western\textquotedblright{}.) Say it once more.
\end{tcolorbox}
\section[il punto di riferimento]{il punto di riferimento — the landmark, the reference point}

\label{lesson:IT-C443-punto-di-riferimento}



\begin{quote}
Before the new one: say the Italian for eastern, then the Italian for western.
\end{quote}

\subsection*{You'll want to know: il punto di riferimento}
\textbf{il punto di riferimento} — \textquotedblleft{}the landmark, the reference point\textquotedblright{}.

\begin{quote}
La torre è un buon punto di riferimento.
\end{quote}

\begin{quote}
\emph{The tower is a good landmark.}
\end{quote}

Also a person you rely on: \emph{mia madre è il mio punto di riferimento}.

\begin{cousinweb}
\emph{punto di riferimento} is built from \emph{punto} 'point', inherited from Latin \emph{punctum} (EtymDB), and \emph{riferimento}, from \emph{riferire}, which comes from Latin \emph{referre} 'to carry back' (EtymDB).
\end{cousinweb}

\begin{grammarlens}[title={What you've built}]
One more word for telling, explaining and describing at B1.
\end{grammarlens}

\subsection*{Your turn}
\begin{itemize}
\raggedright
  \item \emph{Say it:} \emph{il punto di riferimento}
  \item \emph{Say it:} \emph{il punto di riferimento}, once more
  \item \emph{Say it:} a sentence of your own with \emph{il punto di riferimento}
  \item \emph{Recall:} say the Italian for eastern, then the Italian for western, then say \emph{il punto di riferimento} again
\end{itemize}

\begin{tcolorbox}[breakable,colback=teal!4,colframe=teal!35!black,title={Before you move on}]
What does il punto di riferimento mean? (\textquotedblleft{}The landmark, the reference point\textquotedblright{}.) Say it once more.
\end{tcolorbox}
